% =====================================================================
% Section 3: Problem Formulation
% =====================================================================
%
% Status: Draft §3 Problem Formulation v4 (2026-05-19)
%   v4: quarter-page tighten — 100w target.
%       Single ranked-list interface + evidence-set objective in one
%       compact aligned block; fact unit reduced to one inline line.
%       ConnCov instantiation references moved to §4.
%   v3: see archive/03_problem_formulation_v3_pre-quarter-page_20260519.tex
%
% =====================================================================

\section{Problem Formulation}
\label{sec:problem}

Given a question $q$ and a corpus $\mathcal{D}=\{d_i\}_{i=1}^{N}$,
the retriever outputs a ranked evidence list
$R_q=(d_{(1)},d_{(2)},\ldots)$ and the reader consumes its prefix
$R_q^{(K)}=(d_{(1)},\ldots,d_{(K)})$ to produce
$\hat{a}_q=\mathrm{LLM}(\mathrm{prompt}(q,R_q^{(K)}))$. Multi-hop
QA further requires $R_q^{(K)}$ to be \emph{connected} across
documents and \emph{covering} of the entities and relations $q$
asks about. Alongside the ranked list, we therefore frame
retrieval as an evidence-set objective
\begin{equation}
\label{eq:evidence-set}
E_q^{\star}
=
\operatorname*{arg\,max}_{E \subseteq \mathcal{D},\,|E|=K}
\;\mathrm{ConnCov}(E\,;\,q),
\end{equation}
where $\mathrm{ConnCov}$ is instantiated by the query-local graph
linking of \S\ref{sec:method:retrieval} (connectedness) and the
noisy-OR coverage of \S\ref{sec:method:enhancement} (coverage). At
indexing time, each fact unit $f=(h,r,t)$ retains its provenance
$\mathrm{src}(f)\in\mathcal{D}$, and \methodname{} uses these
provenance-linked facts as evidence records for document transitions and rank
passages so that prefixes of $R_q$
approximate $E_q^{\star}$.
